\part{Stand von Forschung und Praxis}
\chapterimage{Bilder_und_Co/autonomer_roboter.jpg} % Chapter heading image
%----------------------------------------------------------------------------------------
% --- KAPITEL: LITERATUR- UND PRAXISÜBERBLICK ---
\chapter{Stand von Forschung \& Praxis}\label{chap:state-of-practice}

Dieses Kapitel skizziert den aktuellen Einsatz von KI – mit klarem Fokus auf die Luftfahrt, jedoch nicht exklusiv –,
zeigt die Schwerpunkte des Forschungsdiskurses zu \emph{AI Safety} und wiederkehrende Praxisprobleme.
Die Darstellung ist bewusst bewertungsfrei und dient als \emph{Context-Setting} zwischen den
\emph{Grundlagen} und den \emph{Anforderungen und Szenarien}.


\section{Grenzen klassischer Verfahren und Motivation}\label{sec:motivation-classic-limits}
Die Luftfahrt ist eine sicherheitskritische Domäne mit dynamischen und komplexen Bedingungen. Bewährte, regel- und modellbasierte Verfahren
bilden den Grundstock, stoßen jedoch sichtbar an Grenzen:

\begin{itemize}
  \item \textbf{Skalierung:} Mit wachsenden Sensor-Datenraten steigen Pflege- und Rechenaufwände überproportional; komplexe Regel-/Filterketten werden zur Engstelle.
  \item \textbf{Nichtlineare Bedingungen:} Fest verdrahtete Regeln bilden stark nichtlineare Zusammenhänge und Wechselwirkungen (z.\,B.\ Geometrie $\times$ Wetter $\times$ Clutter) nur unzureichend ab.
  \item \textbf{Hoher Anpassungsaufwand:} Neue Plattformen/Umgebungen erfordern wiederholtes Nachjustieren von Schwellen/Regeln; Nebenwirkungen sind schwer vorherzusagen.
  \item \textbf{Menschlicher Engpass:} Mit wachsender Datenmenge und -dimensionalität wird die manuelle Musterfindung zum Flaschenhals; Visualisierung allein genügt nicht, es braucht Werkzeuge zur belastbaren Mustererkennung und Verdichtung.
\end{itemize}

\paragraph{Motivation}
KI ist heute allgegenwärtig – von generativen Assistenten wie ChatGPT bis zur KI-gestützten Suchlogik bei jeder Google-Suche; 
hinzu kommen (teil-)autonome Fahrfunktionen und der Wettlauf um vollautonomes Fahren. Der aktuelle „KI-Hype“ schürt hohe 
Erwartungen an Leistungsfähigkeit und Einführungs­tempo. In der Luftfahrt sind die Erwartungen traditionell besonders hoch: 
Entscheidend sind neben Leistungsgewinnen vor allem belastbare Nachweise zur \emph{Sicherheit}. 
KI lernt Muster in hochdimensionalen Daten und adressiert damit \emph{Nichtlinearität} und \emph{Adaptivität} 
(z.\,B.\ Re-Training statt Neuentwurf). Im Radar-Kontext kann KI Muster \emph{im} Clutter erkennen, 
während klassische Verfahren Clutter primär unterdrücken und schwache Ziele ggf.\ erst später sichtbar werden.


% ----------------------------------------------------
\section{KI in der Praxis}\label{sec:praxis}
KI ist längst im industriellen Alltag angekommen – und 2025 zeichnet sich als Beschleunigerjahr ab.
Auf Leitmessen wie der Hannover Messe zeigen Hersteller, wie KI in Serienprozesse wandert.
Laut einer Bitkom-Befragung setzen bereits rund 42\,\% der deutschen Industrieunternehmen KI in der Produktion ein,
weitere 35\,\% planen den baldigen Einsatz.\parencite{handelsblatt-hannovermesse-2025}

\subsection{KI Praxis in der Luftfahrt}
\begin{itemize}
  \item \textbf{Predictive Maintenance:\parencite{handelsblatt-hannovermesse-2025}}
    Rolls-Royce plant „Two-Way-Comms“, um Telemetrie bis auf Komponentenebene aus Triebwerken zu erfassen.
    KI (u.\,a. Microsoft Copilot) soll Anomalien frühzeitig erkennen; Ingenieur:innen können dann gezielt zusätzliche Sensordaten anfordern.
    Der Marktstart wurde für das erste Halbjahr 2025 angekündigt; Recherchen deuten auf Einordnung in das \emph{Intelligent Engine}-Programm hin,
    konkrete belastbare Hinweise auf den produktiven Einsatz dieser KI waren jedoch nicht auffindbar.
  \item \textbf{Design of a nickel-base superalloy:\parencite{conduit2017-superalloy-nn}}
    Kooperation von Rolls-Royce und der University of Cambridge: Ein Deep-Learning-Ansatz schlägt eine neue
    Nickelbasis-Superlegierung für Triebwerkskomponenten vor. Experimente bestätigten die Vorhersagen; die Legierung erfüllte
    alle Zielkriterien und übertraf insbesondere Oxidationsbeständigkeit und Streckgrenze etablierter Alternativen (z.\,B.\ RR1000).
    Öffentliche Informationen zum realen Einsatz dieser Legierung sind derzeit nicht verfügbar.
  \item \textbf{Eurofighter:\parencite{ecrs-mk1-release}}
    Im Eurofighter wird eine von HENSOLDT entwickelte KI-Objektklassifikation im Radar genannt.
    Aufgrund der militärischen Einordnung sind Umsetzungsdetails und Spezifikationen nicht öffentlich;
    der Fall unterstreicht jedoch das Potenzial KI-basierter Objektklassifikation in Luftfahrt-Radarsystemen.
\end{itemize}

\subsection{Blick über den Tellerrand}
Da KI in der Luftfahrt zwar vorhanden, aber im praktischen Einsatz oft schwer nachzuvollziehen ist, lohnt der Blick in andere Branchen, um belastbare Muster und Lessons Learned abzuleiten.

\begin{itemize}
  \item \textbf{Shell – Predictive Maintenance in großem Maßstab:\parencite{shell-jpt-2022}}
    ML-/KI-Modelle überwachen weltweit große Aggregatflotten (Ventile, Pumpen, Kompressoren) und unterstützen 
    zustandsorientierte Wartung (weniger ungeplante Ausfälle, planbarere Einsätze).
    Mit der Erfolgreichen Überwachung von über 10.000 Ausrüstungsteilen plant Shell die Ausweitung dieses Systems.
  \item \textbf{AlphaFold/AlphaFold~3 – Struktur- \& Interaktionsvorhersage:\parencite{alphafold-nature-2021}}
    Breit als Forschungs- und Design-Tool etabliert; hoher wissenschaftlicher Impact, in der Industrie als Enabler akzeptiert 
    (Validierung bleibt erforderlich; keine „KI-entdeckten“ Medikamente allein durch KI zugelassen).
    \emph{Relevanz:} Verdeutlicht die Rolle von KI als \emph{Entscheidungsunterstützung}; 
    Unsicherheit/Vertrauensmaße und experimentelle Verifikation bleiben Pflicht.
  \item \textbf{IBM Watson for Oncology – Grenzen klinischer Empfehlungssysteme:\parencite{stat-watson-2018}}
    Berichte über unsichere/falsche Therapieempfehlungen führten zu Rücknahmen/Neujustierungen prominenter Partnerschaften.
    \emph{Relevanz:} Warnhinweis zu Domain Shift, unzureichender Validierung und HMI: Ohne robuste Evidenz, klare Einsatzgrenzen 
    (ODD) und Human-in-the-Loop riskiert man Scheinsicherheit.
  \item \textbf{KI-gestützte Bildgenerierung – Urheberrecht \& Datenherkunft:\parencite{jipel-andersen-stabilityai-2024}}
    Kontroverse um Trainingsdaten (urheberrechtlich geschützte Werke), laufende Verfahren und Debatten zur Lizenzierung.
    \emph{Relevanz:} Unterstreicht die Bedeutung von Daten-/Label-Governance, Dokumentation der Herkunft und Nutzungsrechten.
\end{itemize}

\subsection{Sicherheits relevante Praxisvorfälle}\label{sec:praxis-vorfaelle}
Die folgenden, öffentlich diskutierten Vorfälle rücken die Bedeutung von \emph{KI-Sicherheit} in den Vordergrund. 
Insbesondere die Automobilbranche steht als treibender Sektor unter hoher öffentlicher Aufmerksamkeit – auch im Hinblick auf 
\emph{eAI (Embodied Artificial Intelligence)} als operierendes Gesamtsystem.

\begin{itemize}
  \item \textbf{Tesla Autopilot:}
  Wiederholt kursieren Berichte und Videos zu unerwartetem Bremsen (sog.\ \emph{Phantom Braking}) oder Abkommen von der Straße
  im Zusammenhang mit dem Tesla Autopiloten im Netz. Der \emph{Autopilot} ist ein Fahrerassistenzsystem und \emph{kein} 
  wie der Name es vermuten lässt, autonomes Fahrzeug. 
  Verantwortung und Überwachung verbleiben daher beim Fahrer. Erkenntnisse (\emph{Lessons Learned}) betreffen vor allem 
  \emph{HMI} (Mensch–Maschine-Interaktion) und \emph{Automation Bias}/Übervertrauen. 
  \item \textbf{Robotaxis (Cruise, San Francisco, 02.10.2023):}
  Nach einer Kollision zwischen einer Fußgängerin und einem anderen Fahrzeug wurde die Fußgängerin unter ein Robotaxi geschleudert. 
  Das Fahrzeug hielt zunächst an, setzte anschließend jedoch kurz an und schleifte die Person einige Meter mit, bevor es erneut stoppte. 
  Dieser Vorfall war unteranderem eine Inspirationen für das Buch \emph{Embodied AI Safety} \parencite{Koopman:2025:EAI}, 
  das \emph{Lessons Learned} und offene Fragen zur AI Safety aufarbeitet.
\end{itemize}

Diese sicherheitsrelevanten Praxisvorfälle werfen zusätzliche Fragen auf. Deutlich wird insbesondere die rechtliche Komplexität: 
Es ist auch zu klären, wer wann die Verantwortung trägt (Fahrer, Betreiber, Hersteller).

% ----------------------------------------------------
\section{Einsatzfelder in der Luftfahrt}\label{sec:sop-ueberblick}
KI findet in der Luftfahrt entlang der gesamten Wertschöpfungskette Anwendungsmöglichkeiten – von Sensorik und Bordfunktionen 
über den Flugbetrieb bis zum Luftraummanagement. Ein Literaturüberblick zeigt folgende Schwerpunkte 
(indikativ; nach \textcite{ml-meets-aerospace-2024}):

\begin{itemize}
  \item \textbf{Flugmanagement und -betrieb (39\,\%):} Routen-/Treibstoffoptimierung, Turnaround- und Dispositionsplanung, Pünktlichkeitsprognosen.
  \item \textbf{Vorhersagende Instandhaltung (34\,\%):} Zustandsdiagnose, Failure Prediction und Flottenanalysen auf Basis von Sensordaten.
  \item \textbf{Sicherheit und Überwachung (23\,\%):} Anomalieerkennung in Netzen/Prozessen, sicherheitsrelevante Muster in Betriebsdaten, Compliance-Monitoring.
  \item \textbf{Luftverkehrsmanagement (4\,\%):} Sequencing/Staffelung, Slot-/Flow-Optimierung, Bodenbewegungen.
\end{itemize}

\paragraph{Einordnung für Radar/Klassifikation.}
Für diese Arbeit ist besonders die KI-gestützte Auswertung von Radardaten (z.\,B.\ Objektklassifikation, Clutter-Trennung, 
Mikrodoppler) relevant. Da luftfahrtspezifische, öffentlich zugängliche Fachquellen dazu derzeit sehr begrenzt sind,
müssen externe Referenzen herangezogen und mit Blick auf die luftfahrttypischen Randbedingungen übertragen werden.
% ----------------------------------------------------
\section{Forschungsstand zu AI Safety}\label{sec:fs-safety}
Der Diskurs priorisiert \emph{Assurance} gegenüber reiner Benchmark-Optimierung. Wiederkehrende Bausteine (anschlussfähig an
die Anforderungsfelder in Kap.~\ref{chap:anforderungen}) sind:
\begin{itemize}
  \item Daten-/Label-Governance \quad (Lineage, Qualität, Freigaben)
  \item ODD \& Annahmenmanagement \quad (klare Einsatzgrenzen, Unknown/Abstain)
  \item Robustheit \& Verteilungssprünge \quad (RFI/Clutter/Wetter/Aspekt)
  \item Unsicherheit \& Kalibrierung \quad (segmentiertes Reporting je ODD)
  \item Erklärbarkeit \& Auditierbarkeit \quad (reproduzierbare Artefakte)
  \item Toolchain-Assurance \& Reproduzierbarkeit \quad (deterministische Builds)
  \item Lebenszyklus/Änderungen \quad (Monitoring, kontrollierte Releases)
\end{itemize}
\textit{Trend:} Der Maßstab verschiebt sich sichtbar von „Accuracy“ zu \emph{Assurance} (ODD-Schnitte, Reproduzierbarkeit,
klare Freigabekriterien).

% ----------------------------------------------------
\section{KI in (Luftfahrt-)Radarsystemen}\label{sec:ki-radar-stand}
In Radarsystemen erscheint KI derzeit in zwei grundsätzlichen Rollen: 
(1) zur \emph{frontloading} Nutzung in Entwicklung, Test und Evidenzerzeugung (offline) sowie 
(2) als \emph{inline} Funktion mit unmittelbarer Wirkung im Betrieb. 
Zu betonen ist die Anwendung im \emph{Frontloading}.

\paragraph{Typische Aufgabenfelder.}
Detektion und Klassifikation (inkl.\ Mikro-Doppler), Ziel/Clutter-Trennung, einfache Sensorfusion, Qualitätsbewertung von
Radarprodukten, kuratierte Datenauswahl/-bereinigung sowie Ableitung von Tests/Szenarien.

\paragraph{Besonderheiten.}
Seltene Ereignisse (Klassenungleichgewicht), verteilungsabhängige Eingaben (Wetter, Seeclutter, RFI, Aspekt) und strikte
Randbedingungen (z.\,B.\ Latenz/Verfügbarkeit bei Bordfunktionen) prägen Datenerhebung, Auswertung und Integration.

\paragraph{Reporting-Praxis.}
Statt Gesamtkennzahlen setzt sich \emph{segmentiertes} Reporting entlang ODD-Schnitten (z.\,B.\ SNR, Entfernung, Wetter, Clutter)
durch, um Leistung und Grenzen kontextsensitiv zu belegen (vgl.\ Kap.~\ref{chap:szenarien}).


% ----------------------------------------------------
\section{Typische Hürden}\label{sec:sop-huerden}
\begin{itemize}
  \item \textbf{Seltene Ereignisse und Class Imbalance:} begrenzte Abdeckung extremer Szenarien; Risiko von Overfitting/Scheingenauigkeit.
  \item \textbf{Verteilungssprünge (Shift):} Wetter/Clutter/RFI und Aspektabhängigkeiten verschieben Datenverteilungen; Kalibrierung kann erodieren.
  \item \textbf{Traceability \& Reproduzierbarkeit:} Ende-zu-Ende-Verfolgbarkeit (Daten $\rightarrow$ Modell $\rightarrow$ Tests $\rightarrow$ Release) als Voraussetzung für Audits.
  \item \textbf{Drift im Betrieb:} saisonale/plattformbedingte Verschiebungen erfordern Monitoring, Re-Kalibrierung und kontrollierte Re-Trainings mit erneuter Freigabe.
  \item \textbf{Integration/HMI:} Qualitätsflags, Konfidenzen und Unknown/Abstain müssen für Operatoren interpretierbar übertragen werden.
\end{itemize}


\section{Besonderheiten der Luftfahrt für KI}\label{sec:ki-besonderheiten}

\subsection{Datenherausforderungen}\label{subsec:data-challenges}
Die Luftfahrt ist geprägt von seltenen Ereignissen und unausgewogenen Klassen:
\begin{itemize}
  \item \textbf{Seltene kritische Situationen} (z.\,B.\ Near-Miss),
  \item \textbf{starke Größen-/Signaturunterschiede} (z.\,B.\ Airbus~A380 vs.\ Vogel),
  \item \textbf{große Geschwindigkeitsvariationen} (Schweben bis Überschall).
\end{itemize}
Seltene Ereignisse erschweren das Training (Klassenimbalance, Overfitting-Risiko); hohe intra-/interklassen-Variabilität erschwert Generalisierung.
Das ist \emph{assurance-relevant}, weil Evidenz segmentiert (ODD-Schnitte) geführt werden muss.

\subsection{Hardware}\label{subsec:hardware}
Je nach Ansatz sind die Anforderungen höher als bei klassischen Verfahren (mehrere CPU-Kerne, ggf.\ GPU/Accelerator).
Solche Komponenten unterliegen Zulassungsanforderungen; verfügbare, qualifizierte Hardware kann den Integrationspfad (Frontloading vs.\ Inline)
und die Latenzbudgets maßgeblich bestimmen. \\
\emph{Praxisnotiz.} In hochkritischen Inline-Pfaden (DAL A/B) werden mehrere \emph{aktive} CPU-Kerne erst seit wenigen Jahren, durch 
jüngerer Behörden-Guidance häufiger tatsächlich eingesetzt; zuvor wurden sie oft deaktiviert, um Nachweise zu vermeiden. 
GPU-/Accelerator-Einsätze sind in solchen Pfaden weiterhin sehr selten. Sie finden sich daher überwiegend in weniger kritischen Partitionen 
oder im Frontloading.



%----------------------------------------------------
\section{Militärische Perspektive – Hauptunterschiede zur zivilen Luftfahrt}\label{sec:militaer-perspektive}
KI findet in der militärischen Luftfahrt bereits vielfältig Anwendung (u.\,a.\ Lagebilder, Sensorik/Wahrnehmung, Mission Support). 
Im Sinne von \emph{Dual-Use} ergeben sich Überschneidungen, von denen zivile und 
militärische Bereiche profitieren können.

Militärische und zivile Verfahren unterscheiden sich grundlegend in Zielsetzung, Einsatzkontext, Regelwerk und Zertifizierung. 
Zivile Luftfahrt priorisiert Passagiersicherheit, Transparenz und internationale Harmonisierung; militärische Luftfahrt fokussiert 
Missionsfähigkeit, Resilienz und Geheimhaltung.
\noindent
Ein aktueller Überblick zur Standardisierung militärischer KI-Anwendungen fasst zentrale Unterschiede zwischen zivilem und militärischem
Kontext wie folgt zusammen:\parencite{mil-ai-standards-2024}

\begin{table}[h]
\centering
\caption{Zivil vs.\ Militär – verkürzter Vergleich zentraler Aspekte}
\begin{tabular}{p{0.30\textwidth} p{0.30\textwidth} p{0.31\textwidth}}
\toprule
\textbf{Aspekt} & \textbf{Zivile Luftfahrt} & \textbf{Militärische Luftfahrt} \\
\midrule
Ziele \& Kontext & Sicherheit, Zuverlässigkeit, Interoperabilität im regulierten Linien-/Frachtbetrieb & Missionswirkung, Robustheit, Anpassungsfähigkeit unter widrigen Bedingungen \\
Regulatorik \& Standards & EASA/FAA/ICAO; DO-178C, DO-254, DO-160; KI-Guidance ergänzt bestehende Safety-Regeln & EDA/NATO/nationale Ministerien; EMARs, MIL-STD-882E; missions-/fähigkeitsbezogene Ergänzungen \\
Standardisierung & Breite Stakeholder-Beteiligung, öffentlich, iterativ & Begrenzte Stakeholder (Behörden, Auftragnehmer), teils klassifiziert \\
Zertifizierung & Pflicht vor Einsatz; internationale Anerkennung; umfangreiche V\&V in kontrollierten Szenarien & EMAR-basiert und missionsbezogen; V\&V auch in Übungen/realitätsnahen Stresstests \\
Transparenz \& Ethik & Hohe Transparenz, Nachvollziehbarkeit, öffentliche Rechenschaft & Abwägung mit Geheimhaltung; zusätzliche ethische Leitplanken (z.\,B.\ zu Autonomiegraden) \\
Betrieb \& Resilienz & „Fail-safe“/„fail-operational“ nach Hazard; klarer Regelbetrieb & Starke Resilienzanforderungen (Jamming/Täuschung), \emph{graceful degradation} unter Störung \\
\bottomrule
\end{tabular}
\end{table}

\paragraph{Kernaussagen.}
\begin{itemize}
  \item \textbf{Ziele:} Zivil ist sicherheits- und compliance-getrieben (maximale Nachweisbarkeit); militärisch missions- und resilienzgetrieben (Wirksamkeit unter Störung).
  \item \textbf{Regelwerke:} Zivil nutzt global harmonisierte, transparente Standards; militärische Regelungen sind stärker auf Einsatzanforderungen und Geheimschutz ausgerichtet.
  \item \textbf{Nachweise:} Zivil erfolgt die Zertifizierung vor Inbetriebnahme mit international anerkannten Evidenzen; militärisch ergänzen realitätsnahe Stresstests/Übungen die V\&V.
  \item \textbf{Transparenz/Ethik:} Zivil betont öffentliche Vertrauensbildung; militärisch wird Transparenz mit operativer Sicherheit und ethischen Leitplanken austariert.
\end{itemize}

\noindent
\textit{Einordnung für diese Arbeit.} Die militärische Perspektive zeigt potenziell höhere Reifegrade einzelner KI-Funktionen (z.\,B.\ Radar),
ist jedoch aufgrund anderer Ziele/Regelwerke nur eingeschränkt auf zivile Zulassung übertragbar. Für die nachfolgenden Anforderungen
und die Coverage-Bewertung bleiben daher zivile Rahmen (EASA, DO-178C/DO-330) ausschlaggebend; militärische Entwicklungen dienen als
Technologieradar, nicht als Zertifizierungsreferenz.






\section{Chancen und Risiken (Radar/Objektklassifizierung)}\label{sec:ki-chancen-risiken}
KI kann Funktionen und Abläufe ergänzen oder verbessern; zugleich entstehen neue Risiken und Herausforderungen:

\begin{table}[h]
\centering
\caption{Chancen vs.\ Risiken/Herausforderungen (Radar/Objektklassifizierung)}
\begin{tabular}{p{0.43\textwidth} p{0.50\textwidth}}
\toprule
\textbf{Chancen} & \textbf{Risiken und Herausforderungen} \\
\midrule
Bessere Trennung Ziel/Clutter & Scheinsicherheit bei Regen-/See-Clutter \\[0.3em]
Frühe Erkennung (z.\,B.\ UAS) & Verpasste Gefahrenklasse (erhöhte False-Negative-Rate) \\[0.3em]
Robustere Sensorfusion & Fehlerfortpflanzung über Sensorkanäle (Kohärenz, Unsicherheitsbehandlung) \\[0.3em]
Anpassbarkeit via Re-Training (ohne Pipeline-Neuentwurf) & Erneuter Freigabe-/Zulassungsbedarf, Daten-Drift im Betrieb \\[0.3em]
Erkennung neuer Muster & Abhängigkeit von Datenqualität/-repräsentativität \\[0.3em]
Entlastung von Personal & Automation-Bias / Fehlinterpretation von Konfidenzen \\
\bottomrule
\end{tabular}
\end{table}


% ----------------------------------------------------
\section{Implikationen für Anforderungen und Nachweis}\label{sec:sop-implikationen}
Auf Basis der Befunde in § \ref{sec:fs-safety}, § \ref{sec:ki-radar-stand} und § \ref{sec:sop-huerden} gelten:
\begin{enumerate}
  \item ODD-/Szenarioabdeckung zuerst
  \item Konfidenz/Kalibrierung verpflichtend
  \item Robustheits- und OOD-Tests
  \item Toolchain-/Daten-Governance
  \item Lifecycle-Sicht
\end{enumerate}

\medskip
\noindent\textit{Übergang.}
Im nächsten Kapitel werden Szenarien konkretisiert, Anforderungsfelder präzisiert
und mit Erwartungen aus Standards/Leitlinien verknüpft, um sie anschließend (u.\,a.\ entlang ISO/PAS~8800)
systematisch auf Abdeckung und Lücken zu prüfen.


\section{Zwischenfazit}\label{sec:ki-zwischenfazit}
KI ist in der Luftfahrt vielversprechend, aber ihre Einführung ist primär ein \emph{Assurance-Problem}, nicht nur ein Trainingsproblem.
Der Maßstab verschiebt sich von „hoher Genauigkeit“ hin zu „\emph{sicher genug}“: Entscheidend sind klare Einsatzgrenzen (ODD),
segmentierte Evidenzen (Leistung \emph{und} Kalibrierung), robuste Nachweise unter Störung sowie geprüfte Degradationspfade.
Diese Beobachtungen leiten direkt zu den \emph{Anforderungen und Szenarien} über.






%--------------------------------------------------------------------------
%komplett unsauber
%------------------------------------------------------------------------------
\chapter{Automobilbranche: KI \& Sicherheit}

\section{Ziel \& Scope}
Dieses Kapitel ordnet KI-Safety in der Automobilbranche ein und nutzt sie als \emph{Vergleichsfolie} für die Luftfahrt:
Welche Bausteine (ODD, szenariobasierte V\&V, Kalibrierung/Unsicherheit, Daten-/Modell-Governance, Monitoring/Updates)
sind etabliert und wie lassen sie sich strukturiert übertragen? Als Referenzrahmen dienen ISO~26262, ISO~21448 (SOTIF),
ISO/PAS~8800 und ISO/IEC~TR~5469.

\subsection{Begriffe \& Standards – Schnellvergleich}
\begin{center}
\begin{tabular}{p{0.38\textwidth} p{0.25\textwidth} p{0.25\textwidth}}
\hline
\textbf{Aspekt} & \textbf{Luftfahrt} & \textbf{Automotive} \\
\hline
Safety-Level & DAL A–E & ASIL A–D \\
System-Safety & ARP4754A / ARP4761A & ISO 26262 (HARA, Safety Concept) / ISO 21448 (SOTIF) \\
Software-Assurance & DO-178C & ISO 26262 Teil SW \\
Hardware-Assurance & DO-254 & ISO 26262 Teil HW \\
KI-Guidance & EASA AI Level 1/2 (Guidance) & ISO/PAS 8800 \\
\hline
\end{tabular}
\end{center}

\section{Normenüberblick (kompakt)}
\begin{itemize}
  \item ISO~26262 (Funktionale Sicherheit)
  \item ISO~21448/SOTIF (beabsichtigte Funktion)
  \item ISO/IEC~TR~5469 (KI+\!Safety-Brücke)
  \item ISO/PAS~8800 (KI-Ergänzung zu 26262/SOTIF)
\end{itemize}
Sie ersetzen keine luftfahrtspezifischen Prozesse, dienen hier aber als Vergleichsfolie für Assurance-Bausteine.\parencite{iso-26262,iso-21448-sotif,iso-iec-tr-5469-2024,iso-pas-8800-2024}

\section{Crosswalk: Luftfahrt \(\leftrightarrow\) Automotive}
\subsection{Gemeinsamkeiten}
\begin{itemize}
  \item Hazard-basierte Skalierung der Nachweisstrenge (DAL \(\leftrightarrow\) ASIL).
  \item Kein „Assurance-Rabatt“ für KI.
  \item Safety Case/Assurance Case als roter Faden (Claims, Strategie, Evidenz, Annahmen).
  \item ODD/Einsatzgrenzen mit Bedeutung von OOD/Unknown und Fallback/Degradation.
  \item Traceability und szenariobasierte V\&V; Toolchain-Betrachtungen bei ML (Daten/Training/Deployment).
\end{itemize}

\subsection{Unterschiede}
\begin{itemize}
  \item \textbf{Regulatorik:} Luftfahrt (Zulassung vor Betrieb, starke Behördenbindung) vs.\ Automotive (EU Typgenehmigung, USA Self-Certification).
  \item \textbf{Änderungen/Updates:} Luftfahrt: erneute Freigabe erforderlich; Automotive: Rückrufaktionen und Feldbeobachtung möglich.
  \item \textbf{HMI \& Betrieb:} Mehrköpfige Crew, Zeitplan vs.\ Privatfahrten.
  \item \textbf{Umgebungen:} Luftfahrt vergleichsweise kontrollierte ODD; Automotive breit gestreute ODDs (urban/land/Autobahn, Wetter, Verkehrsdichte).
  \item \textbf{Physikalische Größen:} Geschwindigkeit, Masse, Abmaße; Luftfahrt stark variierend, Automotive geringere Varianz.
  \item \textbf{Skalen:} Flotten- und Betriebsstundenordnungen differieren (Mengen-/Zeitmaßstäbe), damit auch Evidenz- und Monitoring-Strategien.
  \item \textbf{Einsatzunterscheidung von KI:} nach Autorität/Automatisierung (EASA KI-Level L1/L2 vs.\ SAE Automatisierungsgrad).
\end{itemize}

\subsection{Bemerkungen}
\begin{itemize}
  \item \textbf{DAL VS. ASIL:} DAL berücksichtigt ausschließlich den Schweregrad; ASIL berücksichtigt zusätzlich die 
  Eintrittswahrscheinlichkeit und die Beherrschbarkeit.
  \item \textbf{KI-Level VS. Automatisierungslevel:} KI-Level unterscheidet die Assuranceschicht, welche angelehnt an den 
  Automatisierungsgrad ist; SAE Automatisierungslevel unterscheiden nach Funktionsschicht.
  \item \textbf{Kostenfunktion (Beispiel Radar):} \[
\min_{\tau}\; J(\tau)
= \underbrace{c_{\mathrm{FA}}\;P_{\mathrm{FA}}(\tau)}_{\text{Kosten falscher Alarm}}
+ \underbrace{c_{\mathrm{FN}}\;P_{\mathrm{M}}(\tau)}_{\text{Kosten False Negative}}
+ \underbrace{c_{\mathrm{sys}}(\tau)}_{\text{Systemkosten}} 
\]
  Zivil priorisiert das Vermeiden von Personenschäden; Im militärischen Kontext muss mit Personenschäden gerechnet werden. 
  Weiter kommt hier die Frage auf, ob man die Fehlerwahrscheilichkeit von Menschen im Vergleich  
  zu einer Maschine berücksichtigen sollte.
  \item \textbf{Fehlereintrittsraten – Größenordnungen (Beispielrechnung):} 
  \begin{table}[h]
\centering
\caption{Beispielhafte Dimensionierung für Sensoren (vereinfachtes Rechenbeispiel)}
  \begin{tabular}{p{0.40\textwidth} p{0.25\textwidth} p{0.25\textwidth}}
\toprule
\textbf{Kriterium} & \textbf{Automobilbranche} & \textbf{Luftfahrtbranche} \\
\midrule
  Betriebsstunden pro Fahrzeug (Lebensdauer) & 5.000  & 150.000 \\
  Anzahl verbauter Sensoren (Flotte) & 100.000.000 & 10.000 \\
  Gesamtbetriebsstunden (Flotte) & 5*$10^{11}$h & 15*$10^{8}$h  \\
  Zielraten (Beispiel) & $<10^{-9}$/h &  $<10^{-9}$/h \\
  Anzahl Ausfälle (Betriebsstunden * Zielrate) & < 500 &  < 1{,}5 \\
\bottomrule
\end{tabular}
\end{table}
  \item \textbf{Hohe Sicherheitsanforderungen einzelner Bauteile:} Auch wenn die Motivation unterschiedlich ist, sind 
    die Anforderungen bezüglich der Sicherheitsmetriken wie z.B. die Fehlereintrittsraten ähnlich. Die Automobilbranche 
    versucht Kosten zu vermeiden durch Redundanz und Vermeidung von  Rückrufaktionen und erhöht daher die Sicherheitsanforderungen 
    für die einzelnen Bauteile. Die Luftfahrt hat allgemein hohe Anforderungen an die Bauteile und erhöht die Sicherheit 
    zusätzlich durch Redundanz. 
\end{itemize}



\section{ISO/PAS 8800 — Kernaussagen für diese Arbeit}
\begin{itemize}
  \item \textbf{ODD/Use-Case-Begrenzung} mit Unknown/Abstain \& Fallback als Pflicht.
  \item \textbf{Daten-/Modell-Governance:} Lineage, Versionierung, Label-Qualität, Reproduzierbarkeit.
  \item \textbf{Unsicherheit/Kalibrierung \& OOD:} segmentiertes Reporting, Freigabekriterien.
  \item \textbf{Szenariobasierte V\&V} je ODD-Segment statt Mittelwert-KPIs.
  \item \textbf{Betrieb/Updates:} Monitoring, Drift-Erkennung, kontrollierte Releases/Rollback.
\end{itemize}
\noindent
Diese Punkte werden im Luftfahrtteil nicht 1:1 übernommen, sondern als \emph{Orientierung}
für die bereits definierten Anforderungsfelder genutzt.
\parencite{iso-pas-8800-2024}


\section{Deep Dive: ISO/PAS 8800}
\subsection{Positionierung \& Anwendungsbereich}
\textit{Kernidee:} Ergänzt 26262/SOTIF; kein „Assurance-Rabatt“ für KI; relevant bei ML in sicherheitskritischen Funktionen.

\subsection{Artefakt-Checkliste für den Safety Case}
\textit{Entwürfe:} ODD-Spezifikation; Data/Model-Lineage; Kalibrierungs- \& OOD-Nachweise; Szenariomatrix \& Abdeckung; Runtime-Monitoring-Konzept; Change-/Rollback-Prozess.

\section{ISO/IEC TR 5469 – Brücke \& Ergänzung}
\textit{Rolle:} Referenzrahmen für KI+Funktionale Sicherheit; klärt Einsatzarten (KI \emph{in} Safety-Funktion, Safety \emph{für} KI, KI \emph{für} Entwicklung) und methodische Bausteine (z.\,B.\ Diversität/Unabhängigkeit von KI-Kanälen). 
\textit{Nutzen:} Ergänzt 8800/26262/SOTIF um Begriffe, Risikofaktoren und prozessuale Leitplanken – gut anschlussfähig an Luftfahrt-Argumentationsmuster.
